% GENERATED by report/figures/make_appendix_tables.py from data/wirings/enigma.txt — do not edit.
\begin{table}[ht]\centering\small
\caption{로터와 그리스 로터의 배선(코어 오프셋 0). 입력 글자 A--Z 아래에 출력 글자를 적었다. 역배선은 반사판 쪽에서 들어오는 오는 길이다. 넘김 글자는 그 글자가 창에 보일 때 다음 키에서 왼쪽 로터를 함께 돌린다.}
\label{tab:app-rotors}
\begin{tabular}{@{}l l l c@{}}\toprule
로터 & 배선 (입력 A--Z) & 역배선 & 넘김\\\midrule
입력 & \texttt{ABCDEFGHIJKLMNOPQRSTUVWXYZ} & \texttt{ABCDEFGHIJKLMNOPQRSTUVWXYZ} & \\\midrule
I & \texttt{EKMFLGDQVZNTOWYHXUSPAIBRCJ} & \texttt{UWYGADFPVZBECKMTHXSLRINQOJ} & Q\\
II & \texttt{AJDKSIRUXBLHWTMCQGZNPYFVOE} & \texttt{AJPCZWRLFBDKOTYUQGENHXMIVS} & E\\
III & \texttt{BDFHJLCPRTXVZNYEIWGAKMUSQO} & \texttt{TAGBPCSDQEUFVNZHYIXJWLRKOM} & V\\
IV & \texttt{ESOVPZJAYQUIRHXLNFTGKDCMWB} & \texttt{HZWVARTNLGUPXQCEJMBSKDYOIF} & J\\
V & \texttt{VZBRGITYUPSDNHLXAWMJQOFECK} & \texttt{QCYLXWENFTZOSMVJUDKGIARPHB} & Z\\
VI & \texttt{JPGVOUMFYQBENHZRDKASXLICTW} & \texttt{SKXQLHCNWARVGMEBJPTYFDZUIO} & ZM\\
VII & \texttt{NZJHGRCXMYSWBOUFAIVLPEKQDT} & \texttt{QMGYVPEDRCWTIANUXFKZOSLHJB} & ZM\\
VIII & \texttt{FKQHTLXOCBJSPDZRAMEWNIUYGV} & \texttt{QJINSAYDVKBFRUHMCPLEWZTGXO} & ZM\\
Beta & \texttt{LEYJVCNIXWPBQMDRTAKZGFUHOS} & \texttt{RLFOBVUXHDSANGYKMPZQWEJICT} & ---\\
Gamma & \texttt{FSOKANUERHMBTIYCWLQPZXVGJD} & \texttt{ELPZHAXJNYDRKFCTSIBMGWQVOU} & ---\\
\bottomrule\end{tabular}
\end{table}

\begin{table}[ht]\centering\small
\caption{반사판의 배선. 반사판은 글자를 둘씩 맞바꾸는 대합(involution)이고 고정점이 없다. 그래서 에니그마는 어떤 글자도 자기 자신으로 암호화하지 않는다(크립 정렬의 근거, \ref{sec:crib}절). 얇은 반사판(B-THIN, C-THIN)은 M4에서 그리스 로터와 함께 쓴다.}
\label{tab:app-reflectors}
\begin{tabular}{@{}l l l@{}}\toprule
반사판 & 배선 (입력 A--Z) & 맞바뀌는 쌍\\\midrule
B & \texttt{YRUHQSLDPXNGOKMIEBFZCWVJAT} & \texttt{\footnotesize AY BR CU DH EQ FS GL IP JX KN MO TZ VW}\\
C & \texttt{FVPJIAOYEDRZXWGCTKUQSBNMHL} & \texttt{\footnotesize AF BV CP DJ EI GO HY KR LZ MX NW QT SU}\\
B-THIN & \texttt{ENKQAUYWJICOPBLMDXZVFTHRGS} & \texttt{\footnotesize AE BN CK DQ FU GY HW IJ LO MP RX SZ TV}\\
C-THIN & \texttt{RDOBJNTKVEHMLFCWZAXGYIPSUQ} & \texttt{\footnotesize AR BD CO EJ FN GT HK IV LM PW QZ SX UY}\\
\bottomrule\end{tabular}
\end{table}
